% =====================================================================
%  Thesis chapters - part 1
%  Chapters: 1 Introduction, 2 Theoretical Background
%
%  SELF-CONTAINED DOCUMENT. Compile from the directory containing this
%  file and references.bib, in this order:
%
%      pdflatex chapters
%      bibtex   chapters
%      pdflatex chapters
%      pdflatex chapters
%
%  (or simply:  latexmk -pdf chapters)
%
%  There is deliberately NO \title, NO \author, NO \date and no title
%  page in this file. Add those to your own front matter if needed.
%
%  If you prefer to include these chapters from a larger document
%  instead, delete everything between \begin{document} and the first
%  \chapter, and also delete the \bibliography block at the end of the
%  file; then % =====================================================================
%  Thesis chapters - part 1
%  Chapters: 1 Introduction, 2 Theoretical Background
%
%  SELF-CONTAINED DOCUMENT. Compile from the directory containing this
%  file and references.bib, in this order:
%
%      pdflatex chapters
%      bibtex   chapters
%      pdflatex chapters
%      pdflatex chapters
%
%  (or simply:  latexmk -pdf chapters)
%
%  There is deliberately NO \title, NO \author, NO \date and no title
%  page in this file. Add those to your own front matter if needed.
%
%  If you prefer to include these chapters from a larger document
%  instead, delete everything between \begin{document} and the first
%  \chapter, and also delete the \bibliography block at the end of the
%  file; then % =====================================================================
%  Thesis chapters - part 1
%  Chapters: 1 Introduction, 2 Theoretical Background
%
%  SELF-CONTAINED DOCUMENT. Compile from the directory containing this
%  file and references.bib, in this order:
%
%      pdflatex chapters
%      bibtex   chapters
%      pdflatex chapters
%      pdflatex chapters
%
%  (or simply:  latexmk -pdf chapters)
%
%  There is deliberately NO \title, NO \author, NO \date and no title
%  page in this file. Add those to your own front matter if needed.
%
%  If you prefer to include these chapters from a larger document
%  instead, delete everything between \begin{document} and the first
%  \chapter, and also delete the \bibliography block at the end of the
%  file; then % =====================================================================
%  Thesis chapters - part 1
%  Chapters: 1 Introduction, 2 Theoretical Background
%
%  SELF-CONTAINED DOCUMENT. Compile from the directory containing this
%  file and references.bib, in this order:
%
%      pdflatex chapters
%      bibtex   chapters
%      pdflatex chapters
%      pdflatex chapters
%
%  (or simply:  latexmk -pdf chapters)
%
%  There is deliberately NO \title, NO \author, NO \date and no title
%  page in this file. Add those to your own front matter if needed.
%
%  If you prefer to include these chapters from a larger document
%  instead, delete everything between \begin{document} and the first
%  \chapter, and also delete the \bibliography block at the end of the
%  file; then \input{chapters.tex} from your main file.
%
%  CITATIONS: entries live in references.bib, next to this file. Entries
%  marked "% VERIFY" were written from memory or from a truncated
%  source and should be checked against the original before submission.
% =====================================================================

\documentclass[11pt,a4paper]{report}

% --- encoding, language, fonts -------------------------------------
\usepackage[T1]{fontenc}
\usepackage[utf8]{inputenc}
\usepackage[english]{babel}
\usepackage{lmodern}
\usepackage{microtype}

% --- page layout ---------------------------------------------------
\usepackage[a4paper,margin=2.8cm]{geometry}
\usepackage{setspace}
\onehalfspacing

% --- text ----------------------------------------------------------
\usepackage{amsmath}
\usepackage{amssymb}
\usepackage{caption}
\usepackage{enumitem}
\usepackage{xcolor}

% --- floats, tables, figures ---------------------------------------
\usepackage{graphicx}
\usepackage{booktabs}
\usepackage{float}

% --- citations and cross-references --------------------------------
% Numeric citations: \cite -> [12], \citep -> [12], \citet -> Author [12].
% Swap "numbers" for "authoryear" above if you prefer author-year.
\usepackage[round,numbers]{natbib}
\usepackage[hidelinks]{hyperref}   % load hyperref LAST

% --- list formatting -----------------------------------------------
\setlist[itemize]{leftmargin=1.2em,itemsep=0.3em,topsep=0.4em}
\setlist[enumerate]{leftmargin=1.4em,itemsep=0.3em,topsep=0.4em}

% --- captions ------------------------------------------------------
\captionsetup{font=small,labelfont=bf}

\begin{document}

\tableofcontents

\chapter{Introduction}
\label{ch:introduction}

\section{Motivation}

Geographic information systems (GIS) have become the standard instrument for
working with spatial data at every scale of public administration, science and
industry. The reason for their success is not only analytical power but
\emph{reproducibility}: an analysis performed in a desktop application can be
recorded as a script and re-executed on new data, on another machine, by
another analyst. For a long time this reproducibility was obtained from the
graphical user interfaces of systems such as GRASS \cite{neteler2012grass} and
QGIS \cite{qgissoftware}, whose Python bindings \cite{pyqgis,grass_tools} let
an analyst convert a click sequence into a re-runnable script after the fact.

That order of work, however, is inverted with respect to how analytical needs
are actually formulated. A practitioner usually knows the analysis in
\emph{intent} long before knowing its implementation: it is clear that the
streams must be buffered by 200 metres, or that average temperature must be
computed per administrative unit, while the exact module name, the spelling of
the parameters and the calling convention remain an implementation detail to be
looked up. The gap between intent and implementation is small in knowledge but
large in friction, and it is paid for on every single analysis.

Large language models (LLMs) appear to dissolve that gap: a model can be asked
for an analysis in ordinary language and can answer with code
\cite{vaswani2017attention,chen2021codex}. Since 2023 this capability has been
repackaged into \emph{agents}: systems in which a language model selects tools,
executes them, observes the outcome and iterates \cite{yao2023react}. The
geospatial community has adopted this direction vigorously, both through
research prototypes \cite{zhang2023geogpt,zhang2024geollmtooluse,singh2024geollmengine}
and through a rapidly growing ecosystem of QGIS plugins that put a chat window
in front of a running desktop session.

The enthusiasm, however, meets a hard constraint. The most capable code models
available in practice are large, and the most convenient way of running them is
as a hosted service. For an analyst handling sensitive spatial data --- land
cadastre, utility networks, health data, defence-related imagery --- neither the
upload of the prompt, nor the projection of the resulting code into a vendor's
cloud, nor the loss of control over which model version produced a given
analysis are acceptable. A large part of the current literature on LLM-assisted
GIS silently assumes a cloud endpoint \cite{mansourian2024chatgeoai,singh2024geollmengine}.
An offline alternative must therefore work with what a workstation can run: a
model of a few billion parameters, a local inference server, and the
documentation that ships with the software being automated. Offline-first
operation and hybrid routing are not mutually exclusive: keeping sensitive
requests local while allowing non-sensitive ones to reach a stronger hosted
model gives the analyst both control and capability.

This thesis addresses that gap. It develops and evaluates \textsc{GrassAgent},
an offline-first agent that translates unconstrained natural-language GIS
requests into executable Python scripts written against the modern
object-oriented \texttt{grass.tools} API of GRASS, verifies those scripts
deterministically against the official GRASS manual, and --- as a second stage
--- executes them in a live headless GRASS session and attempts to repair them
from the resulting error feedback.


\chapter{Theoretical Background}
\label{ch:background}

This chapter establishes the conceptual and technical vocabulary on which the
remainder of the thesis depends, and reviews the literature into which the
presented system is placed. Section~\ref{sec:bg:substrate} condenses the
technical substrate: the GIS environment that is automated, the language-model
machinery that performs the translation, the retrieval and routing layers, and
the program-analysis techniques that supply the guarantees.
Section~\ref{sec:bg:related} reviews the prior art, beginning with the method
of the search, and Section~\ref{sec:bg:assistants} examines the operational
assistants most closely comparable to the presented system.

\section{Technical Substrate}
\label{sec:bg:substrate}

\paragraph{The GIS substrate.} GRASS \cite{neteler2012grass} is an open-source
GIS whose 8.6 release exposes more than five hundred and forty modules spanning
raster, vector, imagery, database, display and 3D processing; an agent choosing
among them is solving a selection problem, not a memorisation problem. It is
dual-interface: every module is reachable from the command line and, since
version 8.5, through a machine-generated Python wrapper in which
\texttt{r.slope.aspect} becomes \texttt{r\_slope\_aspect} and every parameter
becomes a keyword argument \cite{grass_tools}. The two interfaces are
functionally equivalent but textually dissimilar, so a model pretrained mostly
on the command line is systematically biased towards the wrong surface syntax.
Execution is also unusually clean --- analysis operates inside a
\emph{location} containing \emph{mapsets}, and every operation is a module
invocation with named options --- which makes process-level isolation of
generated scripts straightforward. The data model distinguishes raster and
vector data, with attribute tables attached to vector features; the distinction
determines which modules may consume which input and is a recurring source of
ambiguity in requests that mix ``each cell'' with ``points''. Add-ons installed
with \texttt{g.extension} register alongside core modules and are reachable
through the same interfaces. QGIS \cite{qgissoftware}, by contrast, exposes
Python bindings \cite{pyqgis} and a Processing framework of several hundred
algorithms invoked with a uniform identifier and parameter dictionary; that
uniform signature explains why much academic work targets PyQGIS
\cite{mansourian2024chatgeoai}, whereas GRASS is comparatively
under-represented and its dual interface is a sharper test of grounding.

\paragraph{Language models for code.} Current code-capable models rest on the
transformer architecture \cite{vaswani2017attention}, and scaling it produced
general-purpose models with emergent code-completion ability
\cite{brown2020gpt3}. Code-specialised models trained on function--documentation
pairs outperform general models of the same size on API-shaped generation
\cite{chen2021codex}, and code-specialised small models such as the
Qwen2.5-Coder family \cite{ji2024qwen25coder} make offline deployment realistic
rather than merely desirable. On a niche API these models hallucinate
identifiers \cite{zan2023surveycode,ji2023hallucination} through three
interacting mechanisms: \emph{surface-form dominance}, since the older and more
frequent command-line form dominates the prior over ``how a GRASS operation
looks in Python''; \emph{compositional pressure}, which tempts the model to
append a plausible but invented finalisation step; and \emph{argument-position
convention transfer}, since the command line names outputs while the wrapper
returns them. These are grounding deficiencies amenable to context and
verification, not reasoning deficiencies that a larger model alone would fix.

\paragraph{Retrieval-augmented generation.} RAG \cite{lewis2020rag} conditions a
parametric model on text fetched from a non-parametric store, addressing
knowledge that is absent from the weights and weights that cannot be updated. For an API surface the store is small and
curated, so retrieval is closer to \emph{interface lookup} than to open-domain
question answering, and precision matters more than coverage. Dense retrieval
\cite{karpukhin2020dpr} is robust to paraphrase; lexical retrieval such as BM25
\cite{robertson2009bm25} is precise when vocabulary matches and brittle when it
does not. On short GIS prompts dense similarity alone tends to promote generic
terrain modules over the specific module named, because a short query embedding
is dominated by the broad topic; hybrid retrieval combines both, and Reciprocal
Rank Fusion \cite{cormack2009rrf} combines rankings by position so that
incomparable score scales do not matter. Domain re-ranking adds priors invisible
to a similarity model: \emph{category weighting} towards raster or vector modules
according to vocabulary, \emph{wrapper preference} for modern wrappers such as
\texttt{v.import} over the raw drivers they replace, and \emph{anti-crowding},
which collapses module families so that siblings cannot crowd out the correct
module. Retrieved context must also be optimised: generic parameters such as
\texttt{flags}, \texttt{overwrite} and \texttt{verbose} cause \emph{parameter
bleed} when left in rendered signatures, and corrupt documentation examples are
faithfully imitated, so deterministic filtering is preferable to instructing the
model to be careful.

\paragraph{Local inference and hybrid routing.} A workstation-class model trades
capability for control: it fits in memory and is fast enough for interactive
use, but is weaker than a hosted frontier model on unfamiliar APIs. Routing
chooses between a local and a remote backend per request on grounds of data
sensitivity, task difficulty, cost or latency, and local servers such as Ollama
and llama.cpp expose an OpenAI-compatible interface so that the same client can
address either. A verification-first design tolerates a weaker local model
because the validator rejects malformed output regardless of its origin, while
cloud routing is acceptable only when the prompt contains no data that may not
leave the host.

\paragraph{Verification and program analysis.} Static verification over the
Python abstract syntax tree \cite{astech} checks each generated call against a
schema for module existence, required parameters, unknown parameters, the
keyword-only calling convention, and dictionary unpacking, which has no name in
the tree and would otherwise escape the check. It is exact with respect to its
schema and reproducible --- no sampling, no temperature, no model in the loop
--- but says nothing about semantic correctness or whether input layers exist.
The complementary technique executes the artifact and returns the observation
as evidence. Automated program repair either patches a failure location
\cite{gazzola2019automatic} or re-prompts with the failing program and its
output, and Reflexion \cite{shinn2023reflexion} established the general form of
the latter as verbal reinforcement, placing the model's own failure observation
back into context without weight updates. Runtime feedback catches failures that
static analysis cannot see --- a missing layer, a type error, a wrong column
name --- at the risk of unbounded cost and non-convergence, which calls for an
attempt ceiling and recorded failure modes.

\section{Related Work: Natural Language to GIS}
\label{sec:bg:related}

The review below was conducted as a structured search rather than as an
accumulation of reading. Works were included if they translate natural language
into geospatial code or workflows, evaluate language models on geospatial code,
present an operational language-model-driven GIS assistant, or supply the
generic machinery --- retrieval, agents, program analysis --- on which such a
system depends; purely conversational assistants and remote-sensing
vision-language work were excluded as non-comparable. Searches covered the
academic indexes and preprint servers, supplemented by citation chasing and by
a systematic survey of the official QGIS plugin repository. Because the field
and the plugin ecosystem move quickly, comparative statements reflect the state
at the date of access.

\subsection{The Autonomous GIS Vision}

The earliest substantial statement of the problem addressed here is the
\emph{Autonomous GIS} proposal, in which a language model acts as the reasoning
engine of a GIS, decomposing a task into sub-tasks and generating code for each
\cite{li2023autonomousgis}. The contribution was primarily conceptual: it
established that the interesting target is not a single function call but a
workflow, and that the bottleneck is task decomposition rather than syntax.

\subsection{Grounded, Tool-Augmented Assistants}

GeoGPT extended a general-purpose language model with a set of tools for
spatial queries and analysis, and demonstrated end-to-end GIS workflows driven
by natural language \cite{zhang2023geogpt}. Subsequent work along the same line
developed a fine-tuned variant operating GIS tools from annotated
instruction--solution data \cite{wei2024geotoolgpt}, and a smaller model adapted
to GIS question answering from domain corpora \cite{zhang2024bbgeogpt}. The
characteristic of this family is that grounding is achieved by \emph{tool
provisioning} --- the agent is given a catalogue and asked to select --- rather
than by verification of the artifact produced.

A parallel line approached the problem from the generation side, training a
language model specifically to emit tool-use chains for geospatial analysis in
a simulated environment \cite{zhang2024geollmtooluse}, and to generate
executable PyQGIS code \cite{mansourian2024chatgeoai}. This line makes the
target unambiguous: the unit of output is a script, and the
evaluation is therefore about code. It also illustrates the dependency of the
whole line on a fine-tuning corpus that, in the published instances, is
proprietary in its preparation and expensive to reproduce --- a property that
motivates a retrieval-based alternative requiring no training data.

\subsection{Benchmarking and Evaluation Practice}

Evaluation of these systems has matured faster than the systems themselves, and
the benchmark literature is arguably the most useful contribution of the recent
period to anyone building such a tool.

A dedicated evaluation of code language models on geospatial code generation
found that performance varies sharply by task framing, and that models
plausible under loose evaluation can fail under strict signature-level checking
\cite{evalcodellms2024}. This is the finding that motivates the design of the
benchmark used in this thesis, whose pass criteria are deliberately strict and
whose verdict depends on the exact interface rather than on plausibility.

GeoLLM-Engine constructed an environment of more than a million queries over
long-horizon geospatial tasks, and reported success rates near 77\% for a
commercial agent \cite{singh2024geollmengine}. Its contribution is the
recognition that single-task, template-based benchmarks are unrepresentative:
real analytical requests are long-horizon and require correct tool sequences.
GeoAnalystBench took the complementary direction, assembling a benchmark of
spatial analysis workflows with human-in-the-loop annotation to construct
reference solutions \cite{zhang2025geoanalystbench}.

The gap these benchmarks expose is the one this thesis occupies. Benchmarks
exist; offline, reproducible generation systems that can be evaluated against
them with static guarantees do not.

\subsection{Multi-Agent and Hybrid Architectures}

A more recent line combines chain-of-thought reasoning, retrieval and
specialised collaborating agents, with fine-tuning to improve spatial task
comprehension and tool selection \cite{mansourian2026multia}. The direction is
promising and the reported results encouraging, but the architecture assumes a
hosted multi-agent service, which places it outside the offline constraint that
governs this work.

Table~\ref{tab:related} summarises the position of the reviewed systems along
the dimensions that distinguish them for the present purpose.

\begin{table}[htbp]
\centering
\caption{Comparison of reviewed natural-language-to-GIS systems. ``Grounding''
denotes the mechanism by which the model is constrained to the real API;
``Verification'' denotes whether the artifact is checked before or after
generation.}
\label{tab:related}
\small
\begin{tabular}{p{2.6cm} p{1.9cm} p{3.1cm} p{2.4cm} p{2.6cm}}
\toprule
\textbf{System} & \textbf{Platform} & \textbf{Grounding} & \textbf{Execution} &
\textbf{Verification} \\
\midrule
GeoGPT \cite{zhang2023geogpt} & Platform-agnostic & Tool catalogue & Agent-internal &
Tool-level \\
Autonomous GIS \cite{li2023autonomousgis} & Platform-agnostic & Task decomposition
& Generated + executed & None reported \\
GeoTool-GPT \cite{wei2024geotoolgpt} & Commercial GIS API & Fine-tuning & Tool use
via LLM & None \\
ChatGeoAI \cite{mansourian2024chatgeoai} & QGIS / PyQGIS & Fine-tuning + RAG
& Script generation & None \\
GeoLLM-Engine \cite{singh2024geollmengine} & Earth observation & Tool catalogue
& Engine-internal & Model-checkers \\
GeoPilot & QGIS & Project context (live layers, CRS, fields) &
\texttt{exec()} in QGIS & Regex guardlist, user confirmation \\
GeoAgent & QGIS, leafmap, anymap & Structured tool registry & Tool calls +
gated script fallback & Confirmation hooks, metadata flags \\
\midrule
\textsc{GrassAgent} (this work) & GRASS & Retrieval over
documentation-derived schema & Headless subprocess, isolated mapset & Static
AST + schema; runtime classification \\
\bottomrule
\end{tabular}
\end{table}

\section{Operational Assistants and Their Design Choices}
\label{sec:bg:assistants}

Two operational systems illustrate the design space most directly:
\textsc{GeoPilot}, a QGIS plugin, and \textsc{GeoAgent}, a library with a QGIS
front-end. Neither is an integration target for this thesis, but their choices
clarify what a benchmarkable, offline-first system must do differently.

\paragraph{\textsc{GeoPilot}.} \textsc{GeoPilot} \cite{geopilot} lets a user
drive a live QGIS session from a chat panel. It gathers the project's layers,
coordinate reference systems, fields and available Processing algorithms,
prompts a model with that state, parses fenced code blocks, and executes the
result with \texttt{exec()} in a namespace containing QGIS objects. By default
the user confirms the code before it runs; an optional auto-execute mode bypasses
the confirmation, and a regular-expression guard blocks a short list of dangerous
constructs (\texttt{os.system}, \texttt{subprocess}, \texttt{shutil.rmtree},
\texttt{eval}). Model access is broad and provider-agnostic, including a local
Ollama backend.

Its grounding mechanism is \emph{runtime project state}: the model is told what
exists in the project rather than what the API permits. That is well suited to
interactive work, where the user's layers are the usual source of error, but it
leaves interface correctness entirely to the model. There is no signature-level
check, and a code error surfaces as a traceback for the user to repair. Execution
is unsandboxed at the Python level, mitigated only by confirmation, and the
plugin offers no headless, reproducible evaluation path.

\paragraph{\textsc{GeoAgent}.} \textsc{GeoAgent} \cite{geoagent} is a shared
agent layer, built on the Strands Agents framework, that exposes geospatial
Python functionality to a model through a structured tool registry. Its QGIS
interface is distributed as the \textsc{OpenGeoAgent} plugin; the same library
also binds to \texttt{leafmap}, \texttt{anymap}, SpatioTemporal Asset
Catalogues and NASA Earthdata. Four abstractions organise it: a
\texttt{GeoAgent} facade, \texttt{GeoAgentConfig} for provider and sampling
settings, \texttt{GeoAgentContext} for live runtime objects, and a
\texttt{@geo\_tool} decorator that turns a typed function into a tool with its
signature and docstring. Tools are grouped into inspection, navigation, layer
management, data loading, selection and processing, and a confirmation-gated
PyQGIS script fallback covers tasks without a dedicated tool. Tools carry risk
metadata such as \texttt{requires\_confirmation}, \texttt{destructive} and
\texttt{long\_running}, which hooks intercept before execution.

Two choices matter here. First, \emph{structural grounding}: the tool surface is
declared statically, with typed signatures and docstrings, which is exactly the
machine-checkable interface representation that deterministic verification
needs. Second, \emph{risk metadata as a first-class property}, which generalises
the question of what generated code is permitted to do. The limitations are the
mirror image: the integration is inherently interactive (a bound map, a QGIS
interface, a dockable panel), it documents no headless evaluation mode, and a
task outside the hand-curated tool surface must fall back on unverified
generated code.

\paragraph{Other implementations.} A survey of the QGIS plugin repository shows
that the pattern is widespread. Several plugins bundle catalogues of tool and
API definitions, retrieve from them before execution, and archive successful
invocations; others advertise very large tool surfaces backed by a commercial
model, or delegate entirely to aggregator APIs. What distinguishes the two
systems above is not membership in this family but the articulation of their
design: a documented grounding mechanism, an explicit execution model, and an
explicit safety posture.

\paragraph{Lessons for the present work.} \textsc{GrassAgent} adopts structural
grounding and risk metadata, but derives the interface from the official manual
rather than from hand-written declarations, and replaces interactive
confirmation with process isolation because its evaluation is unattended. The
comparison is therefore complementary rather than competitive: the two systems
optimise for a present user in a live session, while this work optimises for a
reproducible batch guarantee.

% ---------------------------------------------------------------------
%  Bibliography. Delete this block if these chapters are \input into a
%  larger document that prints its own bibliography.
% ---------------------------------------------------------------------
\bibliographystyle{plainnat}
\bibliography{references}

\end{document} from your main file.
%
%  CITATIONS: entries live in references.bib, next to this file. Entries
%  marked "% VERIFY" were written from memory or from a truncated
%  source and should be checked against the original before submission.
% =====================================================================

\documentclass[11pt,a4paper]{report}

% --- encoding, language, fonts -------------------------------------
\usepackage[T1]{fontenc}
\usepackage[utf8]{inputenc}
\usepackage[english]{babel}
\usepackage{lmodern}
\usepackage{microtype}

% --- page layout ---------------------------------------------------
\usepackage[a4paper,margin=2.8cm]{geometry}
\usepackage{setspace}
\onehalfspacing

% --- text ----------------------------------------------------------
\usepackage{amsmath}
\usepackage{amssymb}
\usepackage{caption}
\usepackage{enumitem}
\usepackage{xcolor}

% --- floats, tables, figures ---------------------------------------
\usepackage{graphicx}
\usepackage{booktabs}
\usepackage{float}

% --- citations and cross-references --------------------------------
% Numeric citations: \cite -> [12], \citep -> [12], \citet -> Author [12].
% Swap "numbers" for "authoryear" above if you prefer author-year.
\usepackage[round,numbers]{natbib}
\usepackage[hidelinks]{hyperref}   % load hyperref LAST

% --- list formatting -----------------------------------------------
\setlist[itemize]{leftmargin=1.2em,itemsep=0.3em,topsep=0.4em}
\setlist[enumerate]{leftmargin=1.4em,itemsep=0.3em,topsep=0.4em}

% --- captions ------------------------------------------------------
\captionsetup{font=small,labelfont=bf}

\begin{document}

\tableofcontents

\chapter{Introduction}
\label{ch:introduction}

\section{Motivation}

Geographic information systems (GIS) have become the standard instrument for
working with spatial data at every scale of public administration, science and
industry. The reason for their success is not only analytical power but
\emph{reproducibility}: an analysis performed in a desktop application can be
recorded as a script and re-executed on new data, on another machine, by
another analyst. For a long time this reproducibility was obtained from the
graphical user interfaces of systems such as GRASS \cite{neteler2012grass} and
QGIS \cite{qgissoftware}, whose Python bindings \cite{pyqgis,grass_tools} let
an analyst convert a click sequence into a re-runnable script after the fact.

That order of work, however, is inverted with respect to how analytical needs
are actually formulated. A practitioner usually knows the analysis in
\emph{intent} long before knowing its implementation: it is clear that the
streams must be buffered by 200 metres, or that average temperature must be
computed per administrative unit, while the exact module name, the spelling of
the parameters and the calling convention remain an implementation detail to be
looked up. The gap between intent and implementation is small in knowledge but
large in friction, and it is paid for on every single analysis.

Large language models (LLMs) appear to dissolve that gap: a model can be asked
for an analysis in ordinary language and can answer with code
\cite{vaswani2017attention,chen2021codex}. Since 2023 this capability has been
repackaged into \emph{agents}: systems in which a language model selects tools,
executes them, observes the outcome and iterates \cite{yao2023react}. The
geospatial community has adopted this direction vigorously, both through
research prototypes \cite{zhang2023geogpt,zhang2024geollmtooluse,singh2024geollmengine}
and through a rapidly growing ecosystem of QGIS plugins that put a chat window
in front of a running desktop session.

The enthusiasm, however, meets a hard constraint. The most capable code models
available in practice are large, and the most convenient way of running them is
as a hosted service. For an analyst handling sensitive spatial data --- land
cadastre, utility networks, health data, defence-related imagery --- neither the
upload of the prompt, nor the projection of the resulting code into a vendor's
cloud, nor the loss of control over which model version produced a given
analysis are acceptable. A large part of the current literature on LLM-assisted
GIS silently assumes a cloud endpoint \cite{mansourian2024chatgeoai,singh2024geollmengine}.
An offline alternative must therefore work with what a workstation can run: a
model of a few billion parameters, a local inference server, and the
documentation that ships with the software being automated. Offline-first
operation and hybrid routing are not mutually exclusive: keeping sensitive
requests local while allowing non-sensitive ones to reach a stronger hosted
model gives the analyst both control and capability.

This thesis addresses that gap. It develops and evaluates \textsc{GrassAgent},
an offline-first agent that translates unconstrained natural-language GIS
requests into executable Python scripts written against the modern
object-oriented \texttt{grass.tools} API of GRASS, verifies those scripts
deterministically against the official GRASS manual, and --- as a second stage
--- executes them in a live headless GRASS session and attempts to repair them
from the resulting error feedback.


\chapter{Theoretical Background}
\label{ch:background}

This chapter establishes the conceptual and technical vocabulary on which the
remainder of the thesis depends, and reviews the literature into which the
presented system is placed. Section~\ref{sec:bg:substrate} condenses the
technical substrate: the GIS environment that is automated, the language-model
machinery that performs the translation, the retrieval and routing layers, and
the program-analysis techniques that supply the guarantees.
Section~\ref{sec:bg:related} reviews the prior art, beginning with the method
of the search, and Section~\ref{sec:bg:assistants} examines the operational
assistants most closely comparable to the presented system.

\section{Technical Substrate}
\label{sec:bg:substrate}

\paragraph{The GIS substrate.} GRASS \cite{neteler2012grass} is an open-source
GIS whose 8.6 release exposes more than five hundred and forty modules spanning
raster, vector, imagery, database, display and 3D processing; an agent choosing
among them is solving a selection problem, not a memorisation problem. It is
dual-interface: every module is reachable from the command line and, since
version 8.5, through a machine-generated Python wrapper in which
\texttt{r.slope.aspect} becomes \texttt{r\_slope\_aspect} and every parameter
becomes a keyword argument \cite{grass_tools}. The two interfaces are
functionally equivalent but textually dissimilar, so a model pretrained mostly
on the command line is systematically biased towards the wrong surface syntax.
Execution is also unusually clean --- analysis operates inside a
\emph{location} containing \emph{mapsets}, and every operation is a module
invocation with named options --- which makes process-level isolation of
generated scripts straightforward. The data model distinguishes raster and
vector data, with attribute tables attached to vector features; the distinction
determines which modules may consume which input and is a recurring source of
ambiguity in requests that mix ``each cell'' with ``points''. Add-ons installed
with \texttt{g.extension} register alongside core modules and are reachable
through the same interfaces. QGIS \cite{qgissoftware}, by contrast, exposes
Python bindings \cite{pyqgis} and a Processing framework of several hundred
algorithms invoked with a uniform identifier and parameter dictionary; that
uniform signature explains why much academic work targets PyQGIS
\cite{mansourian2024chatgeoai}, whereas GRASS is comparatively
under-represented and its dual interface is a sharper test of grounding.

\paragraph{Language models for code.} Current code-capable models rest on the
transformer architecture \cite{vaswani2017attention}, and scaling it produced
general-purpose models with emergent code-completion ability
\cite{brown2020gpt3}. Code-specialised models trained on function--documentation
pairs outperform general models of the same size on API-shaped generation
\cite{chen2021codex}, and code-specialised small models such as the
Qwen2.5-Coder family \cite{ji2024qwen25coder} make offline deployment realistic
rather than merely desirable. On a niche API these models hallucinate
identifiers \cite{zan2023surveycode,ji2023hallucination} through three
interacting mechanisms: \emph{surface-form dominance}, since the older and more
frequent command-line form dominates the prior over ``how a GRASS operation
looks in Python''; \emph{compositional pressure}, which tempts the model to
append a plausible but invented finalisation step; and \emph{argument-position
convention transfer}, since the command line names outputs while the wrapper
returns them. These are grounding deficiencies amenable to context and
verification, not reasoning deficiencies that a larger model alone would fix.

\paragraph{Retrieval-augmented generation.} RAG \cite{lewis2020rag} conditions a
parametric model on text fetched from a non-parametric store, addressing
knowledge that is absent from the weights and weights that cannot be updated. For an API surface the store is small and
curated, so retrieval is closer to \emph{interface lookup} than to open-domain
question answering, and precision matters more than coverage. Dense retrieval
\cite{karpukhin2020dpr} is robust to paraphrase; lexical retrieval such as BM25
\cite{robertson2009bm25} is precise when vocabulary matches and brittle when it
does not. On short GIS prompts dense similarity alone tends to promote generic
terrain modules over the specific module named, because a short query embedding
is dominated by the broad topic; hybrid retrieval combines both, and Reciprocal
Rank Fusion \cite{cormack2009rrf} combines rankings by position so that
incomparable score scales do not matter. Domain re-ranking adds priors invisible
to a similarity model: \emph{category weighting} towards raster or vector modules
according to vocabulary, \emph{wrapper preference} for modern wrappers such as
\texttt{v.import} over the raw drivers they replace, and \emph{anti-crowding},
which collapses module families so that siblings cannot crowd out the correct
module. Retrieved context must also be optimised: generic parameters such as
\texttt{flags}, \texttt{overwrite} and \texttt{verbose} cause \emph{parameter
bleed} when left in rendered signatures, and corrupt documentation examples are
faithfully imitated, so deterministic filtering is preferable to instructing the
model to be careful.

\paragraph{Local inference and hybrid routing.} A workstation-class model trades
capability for control: it fits in memory and is fast enough for interactive
use, but is weaker than a hosted frontier model on unfamiliar APIs. Routing
chooses between a local and a remote backend per request on grounds of data
sensitivity, task difficulty, cost or latency, and local servers such as Ollama
and llama.cpp expose an OpenAI-compatible interface so that the same client can
address either. A verification-first design tolerates a weaker local model
because the validator rejects malformed output regardless of its origin, while
cloud routing is acceptable only when the prompt contains no data that may not
leave the host.

\paragraph{Verification and program analysis.} Static verification over the
Python abstract syntax tree \cite{astech} checks each generated call against a
schema for module existence, required parameters, unknown parameters, the
keyword-only calling convention, and dictionary unpacking, which has no name in
the tree and would otherwise escape the check. It is exact with respect to its
schema and reproducible --- no sampling, no temperature, no model in the loop
--- but says nothing about semantic correctness or whether input layers exist.
The complementary technique executes the artifact and returns the observation
as evidence. Automated program repair either patches a failure location
\cite{gazzola2019automatic} or re-prompts with the failing program and its
output, and Reflexion \cite{shinn2023reflexion} established the general form of
the latter as verbal reinforcement, placing the model's own failure observation
back into context without weight updates. Runtime feedback catches failures that
static analysis cannot see --- a missing layer, a type error, a wrong column
name --- at the risk of unbounded cost and non-convergence, which calls for an
attempt ceiling and recorded failure modes.

\section{Related Work: Natural Language to GIS}
\label{sec:bg:related}

The review below was conducted as a structured search rather than as an
accumulation of reading. Works were included if they translate natural language
into geospatial code or workflows, evaluate language models on geospatial code,
present an operational language-model-driven GIS assistant, or supply the
generic machinery --- retrieval, agents, program analysis --- on which such a
system depends; purely conversational assistants and remote-sensing
vision-language work were excluded as non-comparable. Searches covered the
academic indexes and preprint servers, supplemented by citation chasing and by
a systematic survey of the official QGIS plugin repository. Because the field
and the plugin ecosystem move quickly, comparative statements reflect the state
at the date of access.

\subsection{The Autonomous GIS Vision}

The earliest substantial statement of the problem addressed here is the
\emph{Autonomous GIS} proposal, in which a language model acts as the reasoning
engine of a GIS, decomposing a task into sub-tasks and generating code for each
\cite{li2023autonomousgis}. The contribution was primarily conceptual: it
established that the interesting target is not a single function call but a
workflow, and that the bottleneck is task decomposition rather than syntax.

\subsection{Grounded, Tool-Augmented Assistants}

GeoGPT extended a general-purpose language model with a set of tools for
spatial queries and analysis, and demonstrated end-to-end GIS workflows driven
by natural language \cite{zhang2023geogpt}. Subsequent work along the same line
developed a fine-tuned variant operating GIS tools from annotated
instruction--solution data \cite{wei2024geotoolgpt}, and a smaller model adapted
to GIS question answering from domain corpora \cite{zhang2024bbgeogpt}. The
characteristic of this family is that grounding is achieved by \emph{tool
provisioning} --- the agent is given a catalogue and asked to select --- rather
than by verification of the artifact produced.

A parallel line approached the problem from the generation side, training a
language model specifically to emit tool-use chains for geospatial analysis in
a simulated environment \cite{zhang2024geollmtooluse}, and to generate
executable PyQGIS code \cite{mansourian2024chatgeoai}. This line makes the
target unambiguous: the unit of output is a script, and the
evaluation is therefore about code. It also illustrates the dependency of the
whole line on a fine-tuning corpus that, in the published instances, is
proprietary in its preparation and expensive to reproduce --- a property that
motivates a retrieval-based alternative requiring no training data.

\subsection{Benchmarking and Evaluation Practice}

Evaluation of these systems has matured faster than the systems themselves, and
the benchmark literature is arguably the most useful contribution of the recent
period to anyone building such a tool.

A dedicated evaluation of code language models on geospatial code generation
found that performance varies sharply by task framing, and that models
plausible under loose evaluation can fail under strict signature-level checking
\cite{evalcodellms2024}. This is the finding that motivates the design of the
benchmark used in this thesis, whose pass criteria are deliberately strict and
whose verdict depends on the exact interface rather than on plausibility.

GeoLLM-Engine constructed an environment of more than a million queries over
long-horizon geospatial tasks, and reported success rates near 77\% for a
commercial agent \cite{singh2024geollmengine}. Its contribution is the
recognition that single-task, template-based benchmarks are unrepresentative:
real analytical requests are long-horizon and require correct tool sequences.
GeoAnalystBench took the complementary direction, assembling a benchmark of
spatial analysis workflows with human-in-the-loop annotation to construct
reference solutions \cite{zhang2025geoanalystbench}.

The gap these benchmarks expose is the one this thesis occupies. Benchmarks
exist; offline, reproducible generation systems that can be evaluated against
them with static guarantees do not.

\subsection{Multi-Agent and Hybrid Architectures}

A more recent line combines chain-of-thought reasoning, retrieval and
specialised collaborating agents, with fine-tuning to improve spatial task
comprehension and tool selection \cite{mansourian2026multia}. The direction is
promising and the reported results encouraging, but the architecture assumes a
hosted multi-agent service, which places it outside the offline constraint that
governs this work.

Table~\ref{tab:related} summarises the position of the reviewed systems along
the dimensions that distinguish them for the present purpose.

\begin{table}[htbp]
\centering
\caption{Comparison of reviewed natural-language-to-GIS systems. ``Grounding''
denotes the mechanism by which the model is constrained to the real API;
``Verification'' denotes whether the artifact is checked before or after
generation.}
\label{tab:related}
\small
\begin{tabular}{p{2.6cm} p{1.9cm} p{3.1cm} p{2.4cm} p{2.6cm}}
\toprule
\textbf{System} & \textbf{Platform} & \textbf{Grounding} & \textbf{Execution} &
\textbf{Verification} \\
\midrule
GeoGPT \cite{zhang2023geogpt} & Platform-agnostic & Tool catalogue & Agent-internal &
Tool-level \\
Autonomous GIS \cite{li2023autonomousgis} & Platform-agnostic & Task decomposition
& Generated + executed & None reported \\
GeoTool-GPT \cite{wei2024geotoolgpt} & Commercial GIS API & Fine-tuning & Tool use
via LLM & None \\
ChatGeoAI \cite{mansourian2024chatgeoai} & QGIS / PyQGIS & Fine-tuning + RAG
& Script generation & None \\
GeoLLM-Engine \cite{singh2024geollmengine} & Earth observation & Tool catalogue
& Engine-internal & Model-checkers \\
GeoPilot & QGIS & Project context (live layers, CRS, fields) &
\texttt{exec()} in QGIS & Regex guardlist, user confirmation \\
GeoAgent & QGIS, leafmap, anymap & Structured tool registry & Tool calls +
gated script fallback & Confirmation hooks, metadata flags \\
\midrule
\textsc{GrassAgent} (this work) & GRASS & Retrieval over
documentation-derived schema & Headless subprocess, isolated mapset & Static
AST + schema; runtime classification \\
\bottomrule
\end{tabular}
\end{table}

\section{Operational Assistants and Their Design Choices}
\label{sec:bg:assistants}

Two operational systems illustrate the design space most directly:
\textsc{GeoPilot}, a QGIS plugin, and \textsc{GeoAgent}, a library with a QGIS
front-end. Neither is an integration target for this thesis, but their choices
clarify what a benchmarkable, offline-first system must do differently.

\paragraph{\textsc{GeoPilot}.} \textsc{GeoPilot} \cite{geopilot} lets a user
drive a live QGIS session from a chat panel. It gathers the project's layers,
coordinate reference systems, fields and available Processing algorithms,
prompts a model with that state, parses fenced code blocks, and executes the
result with \texttt{exec()} in a namespace containing QGIS objects. By default
the user confirms the code before it runs; an optional auto-execute mode bypasses
the confirmation, and a regular-expression guard blocks a short list of dangerous
constructs (\texttt{os.system}, \texttt{subprocess}, \texttt{shutil.rmtree},
\texttt{eval}). Model access is broad and provider-agnostic, including a local
Ollama backend.

Its grounding mechanism is \emph{runtime project state}: the model is told what
exists in the project rather than what the API permits. That is well suited to
interactive work, where the user's layers are the usual source of error, but it
leaves interface correctness entirely to the model. There is no signature-level
check, and a code error surfaces as a traceback for the user to repair. Execution
is unsandboxed at the Python level, mitigated only by confirmation, and the
plugin offers no headless, reproducible evaluation path.

\paragraph{\textsc{GeoAgent}.} \textsc{GeoAgent} \cite{geoagent} is a shared
agent layer, built on the Strands Agents framework, that exposes geospatial
Python functionality to a model through a structured tool registry. Its QGIS
interface is distributed as the \textsc{OpenGeoAgent} plugin; the same library
also binds to \texttt{leafmap}, \texttt{anymap}, SpatioTemporal Asset
Catalogues and NASA Earthdata. Four abstractions organise it: a
\texttt{GeoAgent} facade, \texttt{GeoAgentConfig} for provider and sampling
settings, \texttt{GeoAgentContext} for live runtime objects, and a
\texttt{@geo\_tool} decorator that turns a typed function into a tool with its
signature and docstring. Tools are grouped into inspection, navigation, layer
management, data loading, selection and processing, and a confirmation-gated
PyQGIS script fallback covers tasks without a dedicated tool. Tools carry risk
metadata such as \texttt{requires\_confirmation}, \texttt{destructive} and
\texttt{long\_running}, which hooks intercept before execution.

Two choices matter here. First, \emph{structural grounding}: the tool surface is
declared statically, with typed signatures and docstrings, which is exactly the
machine-checkable interface representation that deterministic verification
needs. Second, \emph{risk metadata as a first-class property}, which generalises
the question of what generated code is permitted to do. The limitations are the
mirror image: the integration is inherently interactive (a bound map, a QGIS
interface, a dockable panel), it documents no headless evaluation mode, and a
task outside the hand-curated tool surface must fall back on unverified
generated code.

\paragraph{Other implementations.} A survey of the QGIS plugin repository shows
that the pattern is widespread. Several plugins bundle catalogues of tool and
API definitions, retrieve from them before execution, and archive successful
invocations; others advertise very large tool surfaces backed by a commercial
model, or delegate entirely to aggregator APIs. What distinguishes the two
systems above is not membership in this family but the articulation of their
design: a documented grounding mechanism, an explicit execution model, and an
explicit safety posture.

\paragraph{Lessons for the present work.} \textsc{GrassAgent} adopts structural
grounding and risk metadata, but derives the interface from the official manual
rather than from hand-written declarations, and replaces interactive
confirmation with process isolation because its evaluation is unattended. The
comparison is therefore complementary rather than competitive: the two systems
optimise for a present user in a live session, while this work optimises for a
reproducible batch guarantee.

% ---------------------------------------------------------------------
%  Bibliography. Delete this block if these chapters are \input into a
%  larger document that prints its own bibliography.
% ---------------------------------------------------------------------
\bibliographystyle{plainnat}
\bibliography{references}

\end{document} from your main file.
%
%  CITATIONS: entries live in references.bib, next to this file. Entries
%  marked "% VERIFY" were written from memory or from a truncated
%  source and should be checked against the original before submission.
% =====================================================================

\documentclass[11pt,a4paper]{report}

% --- encoding, language, fonts -------------------------------------
\usepackage[T1]{fontenc}
\usepackage[utf8]{inputenc}
\usepackage[english]{babel}
\usepackage{lmodern}
\usepackage{microtype}

% --- page layout ---------------------------------------------------
\usepackage[a4paper,margin=2.8cm]{geometry}
\usepackage{setspace}
\onehalfspacing

% --- text ----------------------------------------------------------
\usepackage{amsmath}
\usepackage{amssymb}
\usepackage{caption}
\usepackage{enumitem}
\usepackage{xcolor}

% --- floats, tables, figures ---------------------------------------
\usepackage{graphicx}
\usepackage{booktabs}
\usepackage{float}

% --- citations and cross-references --------------------------------
% Numeric citations: \cite -> [12], \citep -> [12], \citet -> Author [12].
% Swap "numbers" for "authoryear" above if you prefer author-year.
\usepackage[round,numbers]{natbib}
\usepackage[hidelinks]{hyperref}   % load hyperref LAST

% --- list formatting -----------------------------------------------
\setlist[itemize]{leftmargin=1.2em,itemsep=0.3em,topsep=0.4em}
\setlist[enumerate]{leftmargin=1.4em,itemsep=0.3em,topsep=0.4em}

% --- captions ------------------------------------------------------
\captionsetup{font=small,labelfont=bf}

\begin{document}

\tableofcontents

\chapter{Introduction}
\label{ch:introduction}

\section{Motivation}

Geographic information systems (GIS) have become the standard instrument for
working with spatial data at every scale of public administration, science and
industry. The reason for their success is not only analytical power but
\emph{reproducibility}: an analysis performed in a desktop application can be
recorded as a script and re-executed on new data, on another machine, by
another analyst. For a long time this reproducibility was obtained from the
graphical user interfaces of systems such as GRASS \cite{neteler2012grass} and
QGIS \cite{qgissoftware}, whose Python bindings \cite{pyqgis,grass_tools} let
an analyst convert a click sequence into a re-runnable script after the fact.

That order of work, however, is inverted with respect to how analytical needs
are actually formulated. A practitioner usually knows the analysis in
\emph{intent} long before knowing its implementation: it is clear that the
streams must be buffered by 200 metres, or that average temperature must be
computed per administrative unit, while the exact module name, the spelling of
the parameters and the calling convention remain an implementation detail to be
looked up. The gap between intent and implementation is small in knowledge but
large in friction, and it is paid for on every single analysis.

Large language models (LLMs) appear to dissolve that gap: a model can be asked
for an analysis in ordinary language and can answer with code
\cite{vaswani2017attention,chen2021codex}. Since 2023 this capability has been
repackaged into \emph{agents}: systems in which a language model selects tools,
executes them, observes the outcome and iterates \cite{yao2023react}. The
geospatial community has adopted this direction vigorously, both through
research prototypes \cite{zhang2023geogpt,zhang2024geollmtooluse,singh2024geollmengine}
and through a rapidly growing ecosystem of QGIS plugins that put a chat window
in front of a running desktop session.

The enthusiasm, however, meets a hard constraint. The most capable code models
available in practice are large, and the most convenient way of running them is
as a hosted service. For an analyst handling sensitive spatial data --- land
cadastre, utility networks, health data, defence-related imagery --- neither the
upload of the prompt, nor the projection of the resulting code into a vendor's
cloud, nor the loss of control over which model version produced a given
analysis are acceptable. A large part of the current literature on LLM-assisted
GIS silently assumes a cloud endpoint \cite{mansourian2024chatgeoai,singh2024geollmengine}.
An offline alternative must therefore work with what a workstation can run: a
model of a few billion parameters, a local inference server, and the
documentation that ships with the software being automated. Offline-first
operation and hybrid routing are not mutually exclusive: keeping sensitive
requests local while allowing non-sensitive ones to reach a stronger hosted
model gives the analyst both control and capability.

This thesis addresses that gap. It develops and evaluates \textsc{GrassAgent},
an offline-first agent that translates unconstrained natural-language GIS
requests into executable Python scripts written against the modern
object-oriented \texttt{grass.tools} API of GRASS, verifies those scripts
deterministically against the official GRASS manual, and --- as a second stage
--- executes them in a live headless GRASS session and attempts to repair them
from the resulting error feedback.


\chapter{Theoretical Background}
\label{ch:background}

This chapter establishes the conceptual and technical vocabulary on which the
remainder of the thesis depends, and reviews the literature into which the
presented system is placed. Section~\ref{sec:bg:substrate} condenses the
technical substrate: the GIS environment that is automated, the language-model
machinery that performs the translation, the retrieval and routing layers, and
the program-analysis techniques that supply the guarantees.
Section~\ref{sec:bg:related} reviews the prior art, beginning with the method
of the search, and Section~\ref{sec:bg:assistants} examines the operational
assistants most closely comparable to the presented system.

\section{Technical Substrate}
\label{sec:bg:substrate}

\paragraph{The GIS substrate.} GRASS \cite{neteler2012grass} is an open-source
GIS whose 8.6 release exposes more than five hundred and forty modules spanning
raster, vector, imagery, database, display and 3D processing; an agent choosing
among them is solving a selection problem, not a memorisation problem. It is
dual-interface: every module is reachable from the command line and, since
version 8.5, through a machine-generated Python wrapper in which
\texttt{r.slope.aspect} becomes \texttt{r\_slope\_aspect} and every parameter
becomes a keyword argument \cite{grass_tools}. The two interfaces are
functionally equivalent but textually dissimilar, so a model pretrained mostly
on the command line is systematically biased towards the wrong surface syntax.
Execution is also unusually clean --- analysis operates inside a
\emph{location} containing \emph{mapsets}, and every operation is a module
invocation with named options --- which makes process-level isolation of
generated scripts straightforward. The data model distinguishes raster and
vector data, with attribute tables attached to vector features; the distinction
determines which modules may consume which input and is a recurring source of
ambiguity in requests that mix ``each cell'' with ``points''. Add-ons installed
with \texttt{g.extension} register alongside core modules and are reachable
through the same interfaces. QGIS \cite{qgissoftware}, by contrast, exposes
Python bindings \cite{pyqgis} and a Processing framework of several hundred
algorithms invoked with a uniform identifier and parameter dictionary; that
uniform signature explains why much academic work targets PyQGIS
\cite{mansourian2024chatgeoai}, whereas GRASS is comparatively
under-represented and its dual interface is a sharper test of grounding.

\paragraph{Language models for code.} Current code-capable models rest on the
transformer architecture \cite{vaswani2017attention}, and scaling it produced
general-purpose models with emergent code-completion ability
\cite{brown2020gpt3}. Code-specialised models trained on function--documentation
pairs outperform general models of the same size on API-shaped generation
\cite{chen2021codex}, and code-specialised small models such as the
Qwen2.5-Coder family \cite{ji2024qwen25coder} make offline deployment realistic
rather than merely desirable. On a niche API these models hallucinate
identifiers \cite{zan2023surveycode,ji2023hallucination} through three
interacting mechanisms: \emph{surface-form dominance}, since the older and more
frequent command-line form dominates the prior over ``how a GRASS operation
looks in Python''; \emph{compositional pressure}, which tempts the model to
append a plausible but invented finalisation step; and \emph{argument-position
convention transfer}, since the command line names outputs while the wrapper
returns them. These are grounding deficiencies amenable to context and
verification, not reasoning deficiencies that a larger model alone would fix.

\paragraph{Retrieval-augmented generation.} RAG \cite{lewis2020rag} conditions a
parametric model on text fetched from a non-parametric store, addressing
knowledge that is absent from the weights and weights that cannot be updated. For an API surface the store is small and
curated, so retrieval is closer to \emph{interface lookup} than to open-domain
question answering, and precision matters more than coverage. Dense retrieval
\cite{karpukhin2020dpr} is robust to paraphrase; lexical retrieval such as BM25
\cite{robertson2009bm25} is precise when vocabulary matches and brittle when it
does not. On short GIS prompts dense similarity alone tends to promote generic
terrain modules over the specific module named, because a short query embedding
is dominated by the broad topic; hybrid retrieval combines both, and Reciprocal
Rank Fusion \cite{cormack2009rrf} combines rankings by position so that
incomparable score scales do not matter. Domain re-ranking adds priors invisible
to a similarity model: \emph{category weighting} towards raster or vector modules
according to vocabulary, \emph{wrapper preference} for modern wrappers such as
\texttt{v.import} over the raw drivers they replace, and \emph{anti-crowding},
which collapses module families so that siblings cannot crowd out the correct
module. Retrieved context must also be optimised: generic parameters such as
\texttt{flags}, \texttt{overwrite} and \texttt{verbose} cause \emph{parameter
bleed} when left in rendered signatures, and corrupt documentation examples are
faithfully imitated, so deterministic filtering is preferable to instructing the
model to be careful.

\paragraph{Local inference and hybrid routing.} A workstation-class model trades
capability for control: it fits in memory and is fast enough for interactive
use, but is weaker than a hosted frontier model on unfamiliar APIs. Routing
chooses between a local and a remote backend per request on grounds of data
sensitivity, task difficulty, cost or latency, and local servers such as Ollama
and llama.cpp expose an OpenAI-compatible interface so that the same client can
address either. A verification-first design tolerates a weaker local model
because the validator rejects malformed output regardless of its origin, while
cloud routing is acceptable only when the prompt contains no data that may not
leave the host.

\paragraph{Verification and program analysis.} Static verification over the
Python abstract syntax tree \cite{astech} checks each generated call against a
schema for module existence, required parameters, unknown parameters, the
keyword-only calling convention, and dictionary unpacking, which has no name in
the tree and would otherwise escape the check. It is exact with respect to its
schema and reproducible --- no sampling, no temperature, no model in the loop
--- but says nothing about semantic correctness or whether input layers exist.
The complementary technique executes the artifact and returns the observation
as evidence. Automated program repair either patches a failure location
\cite{gazzola2019automatic} or re-prompts with the failing program and its
output, and Reflexion \cite{shinn2023reflexion} established the general form of
the latter as verbal reinforcement, placing the model's own failure observation
back into context without weight updates. Runtime feedback catches failures that
static analysis cannot see --- a missing layer, a type error, a wrong column
name --- at the risk of unbounded cost and non-convergence, which calls for an
attempt ceiling and recorded failure modes.

\section{Related Work: Natural Language to GIS}
\label{sec:bg:related}

The review below was conducted as a structured search rather than as an
accumulation of reading. Works were included if they translate natural language
into geospatial code or workflows, evaluate language models on geospatial code,
present an operational language-model-driven GIS assistant, or supply the
generic machinery --- retrieval, agents, program analysis --- on which such a
system depends; purely conversational assistants and remote-sensing
vision-language work were excluded as non-comparable. Searches covered the
academic indexes and preprint servers, supplemented by citation chasing and by
a systematic survey of the official QGIS plugin repository. Because the field
and the plugin ecosystem move quickly, comparative statements reflect the state
at the date of access.

\subsection{The Autonomous GIS Vision}

The earliest substantial statement of the problem addressed here is the
\emph{Autonomous GIS} proposal, in which a language model acts as the reasoning
engine of a GIS, decomposing a task into sub-tasks and generating code for each
\cite{li2023autonomousgis}. The contribution was primarily conceptual: it
established that the interesting target is not a single function call but a
workflow, and that the bottleneck is task decomposition rather than syntax.

\subsection{Grounded, Tool-Augmented Assistants}

GeoGPT extended a general-purpose language model with a set of tools for
spatial queries and analysis, and demonstrated end-to-end GIS workflows driven
by natural language \cite{zhang2023geogpt}. Subsequent work along the same line
developed a fine-tuned variant operating GIS tools from annotated
instruction--solution data \cite{wei2024geotoolgpt}, and a smaller model adapted
to GIS question answering from domain corpora \cite{zhang2024bbgeogpt}. The
characteristic of this family is that grounding is achieved by \emph{tool
provisioning} --- the agent is given a catalogue and asked to select --- rather
than by verification of the artifact produced.

A parallel line approached the problem from the generation side, training a
language model specifically to emit tool-use chains for geospatial analysis in
a simulated environment \cite{zhang2024geollmtooluse}, and to generate
executable PyQGIS code \cite{mansourian2024chatgeoai}. This line makes the
target unambiguous: the unit of output is a script, and the
evaluation is therefore about code. It also illustrates the dependency of the
whole line on a fine-tuning corpus that, in the published instances, is
proprietary in its preparation and expensive to reproduce --- a property that
motivates a retrieval-based alternative requiring no training data.

\subsection{Benchmarking and Evaluation Practice}

Evaluation of these systems has matured faster than the systems themselves, and
the benchmark literature is arguably the most useful contribution of the recent
period to anyone building such a tool.

A dedicated evaluation of code language models on geospatial code generation
found that performance varies sharply by task framing, and that models
plausible under loose evaluation can fail under strict signature-level checking
\cite{evalcodellms2024}. This is the finding that motivates the design of the
benchmark used in this thesis, whose pass criteria are deliberately strict and
whose verdict depends on the exact interface rather than on plausibility.

GeoLLM-Engine constructed an environment of more than a million queries over
long-horizon geospatial tasks, and reported success rates near 77\% for a
commercial agent \cite{singh2024geollmengine}. Its contribution is the
recognition that single-task, template-based benchmarks are unrepresentative:
real analytical requests are long-horizon and require correct tool sequences.
GeoAnalystBench took the complementary direction, assembling a benchmark of
spatial analysis workflows with human-in-the-loop annotation to construct
reference solutions \cite{zhang2025geoanalystbench}.

The gap these benchmarks expose is the one this thesis occupies. Benchmarks
exist; offline, reproducible generation systems that can be evaluated against
them with static guarantees do not.

\subsection{Multi-Agent and Hybrid Architectures}

A more recent line combines chain-of-thought reasoning, retrieval and
specialised collaborating agents, with fine-tuning to improve spatial task
comprehension and tool selection \cite{mansourian2026multia}. The direction is
promising and the reported results encouraging, but the architecture assumes a
hosted multi-agent service, which places it outside the offline constraint that
governs this work.

Table~\ref{tab:related} summarises the position of the reviewed systems along
the dimensions that distinguish them for the present purpose.

\begin{table}[htbp]
\centering
\caption{Comparison of reviewed natural-language-to-GIS systems. ``Grounding''
denotes the mechanism by which the model is constrained to the real API;
``Verification'' denotes whether the artifact is checked before or after
generation.}
\label{tab:related}
\small
\begin{tabular}{p{2.6cm} p{1.9cm} p{3.1cm} p{2.4cm} p{2.6cm}}
\toprule
\textbf{System} & \textbf{Platform} & \textbf{Grounding} & \textbf{Execution} &
\textbf{Verification} \\
\midrule
GeoGPT \cite{zhang2023geogpt} & Platform-agnostic & Tool catalogue & Agent-internal &
Tool-level \\
Autonomous GIS \cite{li2023autonomousgis} & Platform-agnostic & Task decomposition
& Generated + executed & None reported \\
GeoTool-GPT \cite{wei2024geotoolgpt} & Commercial GIS API & Fine-tuning & Tool use
via LLM & None \\
ChatGeoAI \cite{mansourian2024chatgeoai} & QGIS / PyQGIS & Fine-tuning + RAG
& Script generation & None \\
GeoLLM-Engine \cite{singh2024geollmengine} & Earth observation & Tool catalogue
& Engine-internal & Model-checkers \\
GeoPilot & QGIS & Project context (live layers, CRS, fields) &
\texttt{exec()} in QGIS & Regex guardlist, user confirmation \\
GeoAgent & QGIS, leafmap, anymap & Structured tool registry & Tool calls +
gated script fallback & Confirmation hooks, metadata flags \\
\midrule
\textsc{GrassAgent} (this work) & GRASS & Retrieval over
documentation-derived schema & Headless subprocess, isolated mapset & Static
AST + schema; runtime classification \\
\bottomrule
\end{tabular}
\end{table}

\section{Operational Assistants and Their Design Choices}
\label{sec:bg:assistants}

Two operational systems illustrate the design space most directly:
\textsc{GeoPilot}, a QGIS plugin, and \textsc{GeoAgent}, a library with a QGIS
front-end. Neither is an integration target for this thesis, but their choices
clarify what a benchmarkable, offline-first system must do differently.

\paragraph{\textsc{GeoPilot}.} \textsc{GeoPilot} \cite{geopilot} lets a user
drive a live QGIS session from a chat panel. It gathers the project's layers,
coordinate reference systems, fields and available Processing algorithms,
prompts a model with that state, parses fenced code blocks, and executes the
result with \texttt{exec()} in a namespace containing QGIS objects. By default
the user confirms the code before it runs; an optional auto-execute mode bypasses
the confirmation, and a regular-expression guard blocks a short list of dangerous
constructs (\texttt{os.system}, \texttt{subprocess}, \texttt{shutil.rmtree},
\texttt{eval}). Model access is broad and provider-agnostic, including a local
Ollama backend.

Its grounding mechanism is \emph{runtime project state}: the model is told what
exists in the project rather than what the API permits. That is well suited to
interactive work, where the user's layers are the usual source of error, but it
leaves interface correctness entirely to the model. There is no signature-level
check, and a code error surfaces as a traceback for the user to repair. Execution
is unsandboxed at the Python level, mitigated only by confirmation, and the
plugin offers no headless, reproducible evaluation path.

\paragraph{\textsc{GeoAgent}.} \textsc{GeoAgent} \cite{geoagent} is a shared
agent layer, built on the Strands Agents framework, that exposes geospatial
Python functionality to a model through a structured tool registry. Its QGIS
interface is distributed as the \textsc{OpenGeoAgent} plugin; the same library
also binds to \texttt{leafmap}, \texttt{anymap}, SpatioTemporal Asset
Catalogues and NASA Earthdata. Four abstractions organise it: a
\texttt{GeoAgent} facade, \texttt{GeoAgentConfig} for provider and sampling
settings, \texttt{GeoAgentContext} for live runtime objects, and a
\texttt{@geo\_tool} decorator that turns a typed function into a tool with its
signature and docstring. Tools are grouped into inspection, navigation, layer
management, data loading, selection and processing, and a confirmation-gated
PyQGIS script fallback covers tasks without a dedicated tool. Tools carry risk
metadata such as \texttt{requires\_confirmation}, \texttt{destructive} and
\texttt{long\_running}, which hooks intercept before execution.

Two choices matter here. First, \emph{structural grounding}: the tool surface is
declared statically, with typed signatures and docstrings, which is exactly the
machine-checkable interface representation that deterministic verification
needs. Second, \emph{risk metadata as a first-class property}, which generalises
the question of what generated code is permitted to do. The limitations are the
mirror image: the integration is inherently interactive (a bound map, a QGIS
interface, a dockable panel), it documents no headless evaluation mode, and a
task outside the hand-curated tool surface must fall back on unverified
generated code.

\paragraph{Other implementations.} A survey of the QGIS plugin repository shows
that the pattern is widespread. Several plugins bundle catalogues of tool and
API definitions, retrieve from them before execution, and archive successful
invocations; others advertise very large tool surfaces backed by a commercial
model, or delegate entirely to aggregator APIs. What distinguishes the two
systems above is not membership in this family but the articulation of their
design: a documented grounding mechanism, an explicit execution model, and an
explicit safety posture.

\paragraph{Lessons for the present work.} \textsc{GrassAgent} adopts structural
grounding and risk metadata, but derives the interface from the official manual
rather than from hand-written declarations, and replaces interactive
confirmation with process isolation because its evaluation is unattended. The
comparison is therefore complementary rather than competitive: the two systems
optimise for a present user in a live session, while this work optimises for a
reproducible batch guarantee.

% ---------------------------------------------------------------------
%  Bibliography. Delete this block if these chapters are \input into a
%  larger document that prints its own bibliography.
% ---------------------------------------------------------------------
\bibliographystyle{plainnat}
\bibliography{references}

\end{document} from your main file.
%
%  CITATIONS: entries live in references.bib, next to this file. Entries
%  marked "% VERIFY" were written from memory or from a truncated
%  source and should be checked against the original before submission.
% =====================================================================

\documentclass[11pt,a4paper]{report}

% --- encoding, language, fonts -------------------------------------
\usepackage[T1]{fontenc}
\usepackage[utf8]{inputenc}
\usepackage[english]{babel}
\usepackage{lmodern}
\usepackage{microtype}

% --- page layout ---------------------------------------------------
\usepackage[a4paper,margin=2.8cm]{geometry}
\usepackage{setspace}
\onehalfspacing

% --- text ----------------------------------------------------------
\usepackage{amsmath}
\usepackage{amssymb}
\usepackage{caption}
\usepackage{enumitem}
\usepackage{xcolor}

% --- floats, tables, figures ---------------------------------------
\usepackage{graphicx}
\usepackage{booktabs}
\usepackage{float}

% --- citations and cross-references --------------------------------
% Numeric citations: \cite -> [12], \citep -> [12], \citet -> Author [12].
% Swap "numbers" for "authoryear" above if you prefer author-year.
\usepackage[round,numbers]{natbib}
\usepackage[hidelinks]{hyperref}   % load hyperref LAST

% --- list formatting -----------------------------------------------
\setlist[itemize]{leftmargin=1.2em,itemsep=0.3em,topsep=0.4em}
\setlist[enumerate]{leftmargin=1.4em,itemsep=0.3em,topsep=0.4em}

% --- captions ------------------------------------------------------
\captionsetup{font=small,labelfont=bf}

\begin{document}

\tableofcontents

\chapter{Introduction}
\label{ch:introduction}

\section{Motivation}

Geographic information systems (GIS) have become the standard instrument for
working with spatial data at every scale of public administration, science and
industry. The reason for their success is not only analytical power but
\emph{reproducibility}: an analysis performed in a desktop application can be
recorded as a script and re-executed on new data, on another machine, by
another analyst. For a long time this reproducibility was obtained from the
graphical user interfaces of systems such as GRASS \cite{neteler2012grass} and
QGIS \cite{qgissoftware}, whose Python bindings \cite{pyqgis,grass_tools} let
an analyst convert a click sequence into a re-runnable script after the fact.

That order of work, however, is inverted with respect to how analytical needs
are actually formulated. A practitioner usually knows the analysis in
\emph{intent} long before knowing its implementation: it is clear that the
streams must be buffered by 200 metres, or that average temperature must be
computed per administrative unit, while the exact module name, the spelling of
the parameters and the calling convention remain an implementation detail to be
looked up. The gap between intent and implementation is small in knowledge but
large in friction, and it is paid for on every single analysis.

Large language models (LLMs) appear to dissolve that gap: a model can be asked
for an analysis in ordinary language and can answer with code
\cite{vaswani2017attention,chen2021codex}. Since 2023 this capability has been
repackaged into \emph{agents}: systems in which a language model selects tools,
executes them, observes the outcome and iterates \cite{yao2023react}. The
geospatial community has adopted this direction vigorously, both through
research prototypes \cite{zhang2023geogpt,zhang2024geollmtooluse,singh2024geollmengine}
and through a rapidly growing ecosystem of QGIS plugins that put a chat window
in front of a running desktop session.

The enthusiasm, however, meets a hard constraint. The most capable code models
available in practice are large, and the most convenient way of running them is
as a hosted service. For an analyst handling sensitive spatial data --- land
cadastre, utility networks, health data, defence-related imagery --- neither the
upload of the prompt, nor the projection of the resulting code into a vendor's
cloud, nor the loss of control over which model version produced a given
analysis are acceptable. A large part of the current literature on LLM-assisted
GIS silently assumes a cloud endpoint \cite{mansourian2024chatgeoai,singh2024geollmengine}.
An offline alternative must therefore work with what a workstation can run: a
model of a few billion parameters, a local inference server, and the
documentation that ships with the software being automated. Offline-first
operation and hybrid routing are not mutually exclusive: keeping sensitive
requests local while allowing non-sensitive ones to reach a stronger hosted
model gives the analyst both control and capability.

This thesis addresses that gap. It develops and evaluates \textsc{GrassAgent},
an offline-first agent that translates unconstrained natural-language GIS
requests into executable Python scripts written against the modern
object-oriented \texttt{grass.tools} API of GRASS, verifies those scripts
deterministically against the official GRASS manual, and --- as a second stage
--- executes them in a live headless GRASS session and attempts to repair them
from the resulting error feedback.


\chapter{Theoretical Background}
\label{ch:background}

This chapter establishes the conceptual and technical vocabulary on which the
remainder of the thesis depends, and reviews the literature into which the
presented system is placed. Section~\ref{sec:bg:substrate} condenses the
technical substrate: the GIS environment that is automated, the language-model
machinery that performs the translation, the retrieval and routing layers, and
the program-analysis techniques that supply the guarantees.
Section~\ref{sec:bg:related} reviews the prior art, beginning with the method
of the search, and Section~\ref{sec:bg:assistants} examines the operational
assistants most closely comparable to the presented system.

\section{Technical Substrate}
\label{sec:bg:substrate}

\paragraph{The GIS substrate.} GRASS \cite{neteler2012grass} is an open-source
GIS whose 8.6 release exposes more than five hundred and forty modules spanning
raster, vector, imagery, database, display and 3D processing; an agent choosing
among them is solving a selection problem, not a memorisation problem. It is
dual-interface: every module is reachable from the command line and, since
version 8.5, through a machine-generated Python wrapper in which
\texttt{r.slope.aspect} becomes \texttt{r\_slope\_aspect} and every parameter
becomes a keyword argument \cite{grass_tools}. The two interfaces are
functionally equivalent but textually dissimilar, so a model pretrained mostly
on the command line is systematically biased towards the wrong surface syntax.
Execution is also unusually clean --- analysis operates inside a
\emph{location} containing \emph{mapsets}, and every operation is a module
invocation with named options --- which makes process-level isolation of
generated scripts straightforward. The data model distinguishes raster and
vector data, with attribute tables attached to vector features; the distinction
determines which modules may consume which input and is a recurring source of
ambiguity in requests that mix ``each cell'' with ``points''. Add-ons installed
with \texttt{g.extension} register alongside core modules and are reachable
through the same interfaces. QGIS \cite{qgissoftware}, by contrast, exposes
Python bindings \cite{pyqgis} and a Processing framework of several hundred
algorithms invoked with a uniform identifier and parameter dictionary; that
uniform signature explains why much academic work targets PyQGIS
\cite{mansourian2024chatgeoai}, whereas GRASS is comparatively
under-represented and its dual interface is a sharper test of grounding.

\paragraph{Language models for code.} Current code-capable models rest on the
transformer architecture \cite{vaswani2017attention}, and scaling it produced
general-purpose models with emergent code-completion ability
\cite{brown2020gpt3}. Code-specialised models trained on function--documentation
pairs outperform general models of the same size on API-shaped generation
\cite{chen2021codex}, and code-specialised small models such as the
Qwen2.5-Coder family \cite{ji2024qwen25coder} make offline deployment realistic
rather than merely desirable. On a niche API these models hallucinate
identifiers \cite{zan2023surveycode,ji2023hallucination} through three
interacting mechanisms: \emph{surface-form dominance}, since the older and more
frequent command-line form dominates the prior over ``how a GRASS operation
looks in Python''; \emph{compositional pressure}, which tempts the model to
append a plausible but invented finalisation step; and \emph{argument-position
convention transfer}, since the command line names outputs while the wrapper
returns them. These are grounding deficiencies amenable to context and
verification, not reasoning deficiencies that a larger model alone would fix.

\paragraph{Retrieval-augmented generation.} RAG \cite{lewis2020rag} conditions a
parametric model on text fetched from a non-parametric store, addressing
knowledge that is absent from the weights and weights that cannot be updated. For an API surface the store is small and
curated, so retrieval is closer to \emph{interface lookup} than to open-domain
question answering, and precision matters more than coverage. Dense retrieval
\cite{karpukhin2020dpr} is robust to paraphrase; lexical retrieval such as BM25
\cite{robertson2009bm25} is precise when vocabulary matches and brittle when it
does not. On short GIS prompts dense similarity alone tends to promote generic
terrain modules over the specific module named, because a short query embedding
is dominated by the broad topic; hybrid retrieval combines both, and Reciprocal
Rank Fusion \cite{cormack2009rrf} combines rankings by position so that
incomparable score scales do not matter. Domain re-ranking adds priors invisible
to a similarity model: \emph{category weighting} towards raster or vector modules
according to vocabulary, \emph{wrapper preference} for modern wrappers such as
\texttt{v.import} over the raw drivers they replace, and \emph{anti-crowding},
which collapses module families so that siblings cannot crowd out the correct
module. Retrieved context must also be optimised: generic parameters such as
\texttt{flags}, \texttt{overwrite} and \texttt{verbose} cause \emph{parameter
bleed} when left in rendered signatures, and corrupt documentation examples are
faithfully imitated, so deterministic filtering is preferable to instructing the
model to be careful.

\paragraph{Local inference and hybrid routing.} A workstation-class model trades
capability for control: it fits in memory and is fast enough for interactive
use, but is weaker than a hosted frontier model on unfamiliar APIs. Routing
chooses between a local and a remote backend per request on grounds of data
sensitivity, task difficulty, cost or latency, and local servers such as Ollama
and llama.cpp expose an OpenAI-compatible interface so that the same client can
address either. A verification-first design tolerates a weaker local model
because the validator rejects malformed output regardless of its origin, while
cloud routing is acceptable only when the prompt contains no data that may not
leave the host.

\paragraph{Verification and program analysis.} Static verification over the
Python abstract syntax tree \cite{astech} checks each generated call against a
schema for module existence, required parameters, unknown parameters, the
keyword-only calling convention, and dictionary unpacking, which has no name in
the tree and would otherwise escape the check. It is exact with respect to its
schema and reproducible --- no sampling, no temperature, no model in the loop
--- but says nothing about semantic correctness or whether input layers exist.
The complementary technique executes the artifact and returns the observation
as evidence. Automated program repair either patches a failure location
\cite{gazzola2019automatic} or re-prompts with the failing program and its
output, and Reflexion \cite{shinn2023reflexion} established the general form of
the latter as verbal reinforcement, placing the model's own failure observation
back into context without weight updates. Runtime feedback catches failures that
static analysis cannot see --- a missing layer, a type error, a wrong column
name --- at the risk of unbounded cost and non-convergence, which calls for an
attempt ceiling and recorded failure modes.

\section{Related Work: Natural Language to GIS}
\label{sec:bg:related}

The review below was conducted as a structured search rather than as an
accumulation of reading. Works were included if they translate natural language
into geospatial code or workflows, evaluate language models on geospatial code,
present an operational language-model-driven GIS assistant, or supply the
generic machinery --- retrieval, agents, program analysis --- on which such a
system depends; purely conversational assistants and remote-sensing
vision-language work were excluded as non-comparable. Searches covered the
academic indexes and preprint servers, supplemented by citation chasing and by
a systematic survey of the official QGIS plugin repository. Because the field
and the plugin ecosystem move quickly, comparative statements reflect the state
at the date of access.

\subsection{The Autonomous GIS Vision}

The earliest substantial statement of the problem addressed here is the
\emph{Autonomous GIS} proposal, in which a language model acts as the reasoning
engine of a GIS, decomposing a task into sub-tasks and generating code for each
\cite{li2023autonomousgis}. The contribution was primarily conceptual: it
established that the interesting target is not a single function call but a
workflow, and that the bottleneck is task decomposition rather than syntax.

\subsection{Grounded, Tool-Augmented Assistants}

GeoGPT extended a general-purpose language model with a set of tools for
spatial queries and analysis, and demonstrated end-to-end GIS workflows driven
by natural language \cite{zhang2023geogpt}. Subsequent work along the same line
developed a fine-tuned variant operating GIS tools from annotated
instruction--solution data \cite{wei2024geotoolgpt}, and a smaller model adapted
to GIS question answering from domain corpora \cite{zhang2024bbgeogpt}. The
characteristic of this family is that grounding is achieved by \emph{tool
provisioning} --- the agent is given a catalogue and asked to select --- rather
than by verification of the artifact produced.

A parallel line approached the problem from the generation side, training a
language model specifically to emit tool-use chains for geospatial analysis in
a simulated environment \cite{zhang2024geollmtooluse}, and to generate
executable PyQGIS code \cite{mansourian2024chatgeoai}. This line makes the
target unambiguous: the unit of output is a script, and the
evaluation is therefore about code. It also illustrates the dependency of the
whole line on a fine-tuning corpus that, in the published instances, is
proprietary in its preparation and expensive to reproduce --- a property that
motivates a retrieval-based alternative requiring no training data.

\subsection{Benchmarking and Evaluation Practice}

Evaluation of these systems has matured faster than the systems themselves, and
the benchmark literature is arguably the most useful contribution of the recent
period to anyone building such a tool.

A dedicated evaluation of code language models on geospatial code generation
found that performance varies sharply by task framing, and that models
plausible under loose evaluation can fail under strict signature-level checking
\cite{evalcodellms2024}. This is the finding that motivates the design of the
benchmark used in this thesis, whose pass criteria are deliberately strict and
whose verdict depends on the exact interface rather than on plausibility.

GeoLLM-Engine constructed an environment of more than a million queries over
long-horizon geospatial tasks, and reported success rates near 77\% for a
commercial agent \cite{singh2024geollmengine}. Its contribution is the
recognition that single-task, template-based benchmarks are unrepresentative:
real analytical requests are long-horizon and require correct tool sequences.
GeoAnalystBench took the complementary direction, assembling a benchmark of
spatial analysis workflows with human-in-the-loop annotation to construct
reference solutions \cite{zhang2025geoanalystbench}.

The gap these benchmarks expose is the one this thesis occupies. Benchmarks
exist; offline, reproducible generation systems that can be evaluated against
them with static guarantees do not.

\subsection{Multi-Agent and Hybrid Architectures}

A more recent line combines chain-of-thought reasoning, retrieval and
specialised collaborating agents, with fine-tuning to improve spatial task
comprehension and tool selection \cite{mansourian2026multia}. The direction is
promising and the reported results encouraging, but the architecture assumes a
hosted multi-agent service, which places it outside the offline constraint that
governs this work.

Table~\ref{tab:related} summarises the position of the reviewed systems along
the dimensions that distinguish them for the present purpose.

\begin{table}[htbp]
\centering
\caption{Comparison of reviewed natural-language-to-GIS systems. ``Grounding''
denotes the mechanism by which the model is constrained to the real API;
``Verification'' denotes whether the artifact is checked before or after
generation.}
\label{tab:related}
\small
\begin{tabular}{p{2.6cm} p{1.9cm} p{3.1cm} p{2.4cm} p{2.6cm}}
\toprule
\textbf{System} & \textbf{Platform} & \textbf{Grounding} & \textbf{Execution} &
\textbf{Verification} \\
\midrule
GeoGPT \cite{zhang2023geogpt} & Platform-agnostic & Tool catalogue & Agent-internal &
Tool-level \\
Autonomous GIS \cite{li2023autonomousgis} & Platform-agnostic & Task decomposition
& Generated + executed & None reported \\
GeoTool-GPT \cite{wei2024geotoolgpt} & Commercial GIS API & Fine-tuning & Tool use
via LLM & None \\
ChatGeoAI \cite{mansourian2024chatgeoai} & QGIS / PyQGIS & Fine-tuning + RAG
& Script generation & None \\
GeoLLM-Engine \cite{singh2024geollmengine} & Earth observation & Tool catalogue
& Engine-internal & Model-checkers \\
GeoPilot & QGIS & Project context (live layers, CRS, fields) &
\texttt{exec()} in QGIS & Regex guardlist, user confirmation \\
GeoAgent & QGIS, leafmap, anymap & Structured tool registry & Tool calls +
gated script fallback & Confirmation hooks, metadata flags \\
\midrule
\textsc{GrassAgent} (this work) & GRASS & Retrieval over
documentation-derived schema & Headless subprocess, isolated mapset & Static
AST + schema; runtime classification \\
\bottomrule
\end{tabular}
\end{table}

\section{Operational Assistants and Their Design Choices}
\label{sec:bg:assistants}

Two operational systems illustrate the design space most directly:
\textsc{GeoPilot}, a QGIS plugin, and \textsc{GeoAgent}, a library with a QGIS
front-end. Neither is an integration target for this thesis, but their choices
clarify what a benchmarkable, offline-first system must do differently.

\paragraph{\textsc{GeoPilot}.} \textsc{GeoPilot} \cite{geopilot} lets a user
drive a live QGIS session from a chat panel. It gathers the project's layers,
coordinate reference systems, fields and available Processing algorithms,
prompts a model with that state, parses fenced code blocks, and executes the
result with \texttt{exec()} in a namespace containing QGIS objects. By default
the user confirms the code before it runs; an optional auto-execute mode bypasses
the confirmation, and a regular-expression guard blocks a short list of dangerous
constructs (\texttt{os.system}, \texttt{subprocess}, \texttt{shutil.rmtree},
\texttt{eval}). Model access is broad and provider-agnostic, including a local
Ollama backend.

Its grounding mechanism is \emph{runtime project state}: the model is told what
exists in the project rather than what the API permits. That is well suited to
interactive work, where the user's layers are the usual source of error, but it
leaves interface correctness entirely to the model. There is no signature-level
check, and a code error surfaces as a traceback for the user to repair. Execution
is unsandboxed at the Python level, mitigated only by confirmation, and the
plugin offers no headless, reproducible evaluation path.

\paragraph{\textsc{GeoAgent}.} \textsc{GeoAgent} \cite{geoagent} is a shared
agent layer, built on the Strands Agents framework, that exposes geospatial
Python functionality to a model through a structured tool registry. Its QGIS
interface is distributed as the \textsc{OpenGeoAgent} plugin; the same library
also binds to \texttt{leafmap}, \texttt{anymap}, SpatioTemporal Asset
Catalogues and NASA Earthdata. Four abstractions organise it: a
\texttt{GeoAgent} facade, \texttt{GeoAgentConfig} for provider and sampling
settings, \texttt{GeoAgentContext} for live runtime objects, and a
\texttt{@geo\_tool} decorator that turns a typed function into a tool with its
signature and docstring. Tools are grouped into inspection, navigation, layer
management, data loading, selection and processing, and a confirmation-gated
PyQGIS script fallback covers tasks without a dedicated tool. Tools carry risk
metadata such as \texttt{requires\_confirmation}, \texttt{destructive} and
\texttt{long\_running}, which hooks intercept before execution.

Two choices matter here. First, \emph{structural grounding}: the tool surface is
declared statically, with typed signatures and docstrings, which is exactly the
machine-checkable interface representation that deterministic verification
needs. Second, \emph{risk metadata as a first-class property}, which generalises
the question of what generated code is permitted to do. The limitations are the
mirror image: the integration is inherently interactive (a bound map, a QGIS
interface, a dockable panel), it documents no headless evaluation mode, and a
task outside the hand-curated tool surface must fall back on unverified
generated code.

\paragraph{Other implementations.} A survey of the QGIS plugin repository shows
that the pattern is widespread. Several plugins bundle catalogues of tool and
API definitions, retrieve from them before execution, and archive successful
invocations; others advertise very large tool surfaces backed by a commercial
model, or delegate entirely to aggregator APIs. What distinguishes the two
systems above is not membership in this family but the articulation of their
design: a documented grounding mechanism, an explicit execution model, and an
explicit safety posture.

\paragraph{Lessons for the present work.} \textsc{GrassAgent} adopts structural
grounding and risk metadata, but derives the interface from the official manual
rather than from hand-written declarations, and replaces interactive
confirmation with process isolation because its evaluation is unattended. The
comparison is therefore complementary rather than competitive: the two systems
optimise for a present user in a live session, while this work optimises for a
reproducible batch guarantee.

% ---------------------------------------------------------------------
%  Bibliography. Delete this block if these chapters are \input into a
%  larger document that prints its own bibliography.
% ---------------------------------------------------------------------
\bibliographystyle{plainnat}
\bibliography{references}

\end{document}